\title{Byte Latent Transformer: Patches Scale Better Than Tokens}

\begin{abstract}
We propose the Byte Latent Transformer ({\textsc BLT}), a novel large language model architecture that operates at the byte level and, for the first time, achieves tokenization-based LLM performance at large scales while also substantially increasing inference robustness and efficiency. The {\textsc BLT} framework maps bytes into variably sized patches, which then become the core computational units of the model. The boundaries of these patches are determined adaptively according to the entropy of upcoming bytes, enabling greater computational resources and model capacity to be directed towards more complex input regions. We conduct the first scaling investigation of byte-level models with strict FLOP budgeting, scaling up to 8B parameters and processing 4T bytes during training. The findings underscore the practicality of training large-scale models directly from byte sequences without relying on a fixed token vocabulary. Dynamic selection of elongated patches in predictable contexts enhances both the training and inference efficiency, alongside observed improvements in advanced reasoning and performance on rare, challenging inputs. Crucially, at equivalent inference cost, {\textsc BLT} demonstrates a much more favorable scaling trajectory relative to token-based models, owing to the concurrent expansion of patch lengths and model size.
\end{abstract}

\section{Introduction}
\label{section:intro}

We present the Byte Latent Transformer~(\textbf{BLT}), a novel architecture that dispenses with tokenization entirely by directly processing sequences of raw bytes. For the first time, our model is able to achieve performance on par with traditional tokenization-based systems at large scales, while offering marked gains in both efficiency and resilience (\S\ref{sec:robustness}). While contemporary large language models (llms) are designed to be nearly fully end-to-end, they rely on tokenization as an initial heuristic phase that converts raw byte sequences into a fixed vocabulary of tokens. This procedure imposes an artificial structure on the input, introducing various issues such as heightened sensitivity to different domains and modalities~\citep{dagangetting}, vulnerability to noisy inputs~(\S\ref{sec:robustness}), inadequate handling of orthographic features~\citep{edman2024cute}, and disadvantages for multilingual processing~\citep{liang2023xlm,petrov2024language,limisiewicz2024myte}.

Tokenization has historically played a crucial role because training llms directly on byte-level inputs incurs substantial computational expense at large scales, mainly due to increased sequence lengths~\citep{xue2022byt5}. Prior research addresses this issue by adopting more efficient self-attention methods~\citep{el2020characterbert, clark2022canine} or leveraging architectures that eliminate attention entirely~\citep{wang2024mambabyte}~(\S\ref{section:related_work}). Nonetheless, these optimizations are primarily advantageous for training \textit{small models}. When working with large models, the main computational bottleneck comes from the extensive feed-forward network layers applied to each byte, rather than from attention mechanisms.

To optimize compute distribution, we introduce a dynamic and learnable approach that organizes bytes into \textit{patches}~(\S\ref{section:patching}), alongside a novel model architecture that integrates both byte-level and patch-level data. Distinct from tokenization, {\textsc BLT} does not rely on a predetermined patch vocabulary. Instead, it utilizes lightweight, trainable encoder and decoder components to transform arbitrary byte groupings into latent patch representations. Our results demonstrate that this mechanism leads to \textit{greater} computational efficiency compared to traditional tokenization-based frameworks.

Tokenization-based language models apply a uniform amount of computational resources to each token, regardless of the specific token’s complexity. This approach favors simplicity over computational efficiency because the heuristics used to define tokens are based on compression strategies rather than the actual difficulty of predicting each token. Our model fundamentally reconsiders this assumption by advocating for adaptive compute allocation in response to the demands of specific prediction tasks. For instance, generating the ending of the majority of words typically poses a lower level of uncertainty and does not require the capabilities of a large-scale transformer—contrasting sharply with the higher entropy involved in selecting the initial word of a sentence. This principle shapes the design of {\textsc BLT}~(\S\ref{section:architecture}), which divides its architecture into three transformer modules: two compact, byte-level \textit{local models}, and a single, much larger, global \textit{latent transformer}~(\autoref{fig:arch_high_level}). 

{\textsc BLT} orchestrates the formation of byte groupings—referred to as patches—by segmenting data in accordance with the entropy associated with the prediction of the upcoming byte. This strategy produces contextualized sets of bytes that exhibit a more even distribution of information content, thereby enabling dynamic and efficient allocation of computational effort.

We introduce the first scaling analysis of byte-level models where training computation is controlled by the number of floating point operations (flops), scaling up to models with 8 billion parameters and training on 4 trillion bytes, and demonstrate that it is feasible to train large-scale models directly on bytes, foregoing the need for fixed-vocabulary tokenization. In our experiments, {\textsc BLT} achieves flop-controlled training performance on par with Llama 3, while delivering up to a 50\% reduction in inference flops (\S\ref{section:scaling}).

Our findings further indicate that models operating directly on raw bytes yield notable gains in capturing the distributional tails of datasets. {\textsc BLT} outperforms traditional tokenizer-based architectures in robustness to input noise, and exhibits superior capabilities on tasks reliant on fine-grained character-level features, including orthographic distinctions, phonological analysis, and low-resource machine translation (\S\ref{sec:robustness}).

Moreover, {\textsc BLT} enables joint increases in both model capacity and byte patch size without exceeding a fixed inference flop allowance. Utilizing larger patch sizes can reduce computational demand per forward pass, freeing resources to allocate towards expanding the global latent transformer component, which operates at lower frequency. Through inference flop-controlled scaling experiments (\autoref{fig:fixed_inference_scaling}), we observe that this approach yields significantly improved scaling behavior compared to traditional tokenization-based frameworks.

To summarize, this work presents several key contributions: 1) We propose {\textsc BLT}, a byte latent LLM framework that adaptively manages computational allocation to enhance flop efficiency; 2) We demonstrate that our method attains flop-matched training performance with Llama 3 up to the 8B parameter range, and can yield flop efficiency improvements up to 50\% by tolerating small reductions in evaluation outcomes; 3) The {\textsc BLT} architecture enables an additional scalability axis for LLMs, allowing expansion in model size without exceeding a fixed inference compute constraint; 4) We provide evidence of greater robustness of {\textsc BLT} models to input perturbations and show that they capture sub-word input features often overlooked by token-based LLMs. All training and inference source code for {\textsc BLT} is made available at~\url{https://github.com/facebookresearch/blt}.

\section{Patching: From Individual Bytes to Groups of Bytes}
\label{section:patching}

Dividing bytes into \textit{patches} enables {\textsc BLT} to flexibly assign computational resources in response to the context. Several approaches to partitioning bytes into patches are illustrated in Figure~\ref{fig:patching_types}. More precisely, a patching function $f_p$ partitions an input sequence of bytes $\pmb{x}=\{x_i,|i=1,\ldots n\}$, of length $n$, into a collection of $m < n$ patches $\pmb{p}=\{p_j|j=1,\ldots,m\}$ by designating for each $x_i$ a mapping into \{0,1\}, where a value of 1 marks the beginning of a new patch. Whether the architecture relies on tokens or patches, the predominant determinant of computational expense lies in how many steps the main Transformer must execute. In {\textsc BLT}, this correlates to the total number of patches generated from the data using a specific patching method. Thus, the average number of bytes per patch—or simply, the \textit{patch size}—serves as the critical measure of computational cost for both training and inference under a chosen patching function~(\S\ref{section:flops}).

In the following, we present three patching strategies: using a fixed byte count per patch~(\S\ref{section:static-patch}), segmenting by whitespace~(\S\ref{section:white-patch}), and adaptively creating patches based on entropy values from a lightweight byte language model~(\S\ref{section:dyn-patch}). Additionally, we cover the method of incremental patching and delineate the distinction between tokenization and patching~(\S\ref{section:bpe-patch}).

\subsection{Strided Patching Every K Bytes}
\label{section:static-patch}

A common method for aggregating bytes involves dividing them into fixed-size segments of length $k$, as utilized in MegaByte~\citep{yu2023megabyte}. Employing a constant stride facilitates both training and inference processes, offers a simple way to adjust the average patch size, and thereby allows for straightforward management of computational cost. Despite its simplicity, this approach is accompanied by notable drawbacks. Firstly, computational resources are allocated uniformly rather than adaptively based on data complexity: for example, a transformer operation $j$ may be inefficiently used when processing low-information content like whitespace in code, or conversely, may provide insufficient capacity for highly informative regions such as mathematical expressions. Secondly, this method can result in inconsistent and context-insensitive segmentation of byte sequences—for instance, identical words may be divided at different points, disrupting their continuity.

\subsection{Space Patching}
\label{section:white-patch}

\citet{slagle2024spacebyte} introduces a straightforward but impactful enhancement to strided patching by segmenting patches after every space-like byte\footnote{
    Space-like bytes refer to any byte that is neither a Latin letter, numeral, nor a UTF-8 continuation byte. Furthermore, patches are required to contain at least one byte that is not space-like. }, which frequently align with linguistic unit boundaries in a wide array of languages. Through Space patching, the model allocates a latent transformer operation (effectively, additional computational resources) to each individual word. This arrangement enables uniform patching of words throughout different sequences and ensures computational attention is directed at tokens that are often challenging to predict—frequently found immediately after spaces. For instance, it is considerably more complex to generate the initial byte of the answer for the prompt ``Who composed the Magic Flute?\uline{\hspace{.6em}}'' compared to completing the rest of the response after the ``M''; once the first character is known, the subsequent bytes are much less ambiguous, so generating ``Mozart'' becomes substantially easier. Nonetheless, space patching displays limitations: it is not universally applicable across all languages and domains, and it does not facilitate flexible patch size adjustment. In the following section, we present a novel patching approach that leverages the observation that initial bytes of words usually present greater prediction difficulty, while also providing a principled means to manage patch size.

\subsection{Entropy Patching: Using Next-Byte Entropies from a Small Byte LM}
\label{section:dyn-patch}

Instead of utilizing a heuristic based on predefined rules like whitespace, we adopt a data-driven strategy to detect locations where next-byte predictions are most uncertain. To this end, we propose \textit{entropy patching}, a method that determines patch boundaries based on calculated entropy values.

We train a small byte-level auto-regressive language model on the training data for {\textsc BLT} and compute next byte entropies under the LM distribution $p_{e}$ over the byte vocabulary $\mathcal{V}$:

\begin{equation}
H(x_i) = \sum_{v \in \mathcal{V}} p_{e}(x_i=v|\pmb{x}_{<i}) \log p_{e}(x_i=v|\pmb{x}_{<i})
\end{equation}

We investigate two approaches for determining patch boundaries using entropies $H(x_i)$. The first method selects points exceeding a predefined global entropy threshold, as depicted in~\autoref{fig:patching}. The second strategy focuses on points where the entropy is substantially higher compared to the preceding value. This latter method effectively seeks points that disrupt the approximately monotonic decline of entropy within a given patch.

\begin{align*}
    \text{Global Constraint}\hspace{1cm}H(x_t)& > \theta_g\\
    \text{Approx. Monotonic Constraint}\hspace{1cm}H(x_t)& - H(x_{t-1}) > \theta_r
\end{align*}

During the dataloading process, a minimal preprocessing step is used to detect patch boundaries. This contrasts with the approach of~\citet{nawrot-etal-2023-efficient}, where a classifier is employed to predict patch boundaries based on entropy measures. In our experiments~(\S\ref{section:experiments}), we evaluate and compare these two techniques for classifying bytes into low and high entropy categories.

\subsection{The Byte-Pair Encoding (BPE) Tokenizer and Incremental Patching}
\label{section:incremental-patching}
\label{section:bpe-patch}

Numerous contemporary llms, such as our baseline Llama 3, rely on subword tokenization methods like bpe~\citep{gage1994new, sennrich-etal-2016-neural}. Throughout this paper, we use the term ``tokens'' to mean groups of bytes selected from a \textit{finite} set that is predefined before training begins, whereas ``patches'' denote dynamically constructed sequences that are not constrained to a predetermined vocabulary. Notably, patches and tokens differ in that tokens prevent the model from accessing the raw byte-level information directly.

An important advancement offered by {\textsc BLT} compared to tokenization-based models is its novel approach to balancing vocabulary size against computational requirements. Conventional large language models (LLMs) face a dilemma: expanding the vocabulary results in, on average, larger tokens and therefore reduces the number of decoding steps needed; however, it simultaneously inflates the output dimensionality of the final projection layer, increasing compute. This inherent compromise restricts the potential of tokenization-based strategies to meaningfully alter token length or reduce inference expenses. As an illustration, Llama 3 achieves a rise in average token size—from 3.7 bytes to 4.4 bytes—but this comes with a fourfold enlargement of its embedding table relative to Llama 2.

During generation, {\textsc BLT} must determine if the current position within the byte sequence corresponds to a patch boundary, since this influences whether the Latent Transformer performs additional computation. This decision must be made without knowledge of future bytes, as those have not been generated yet. Consequently, the patching process must not rely on future information when segmenting the byte sequence. A patching method $f_p$ is said to be incrementally patching if it adheres to the condition:
$$f_p(\pmb{x}_{<i}) = f_p(\pmb{x})_{<i}$$
The bpe method does not fulfill the requirement for incremental patching because identical prefixes can produce different segmentations based on how the sequence is continued, violating the condition above\footnote{Introducing a special delimiter token at patch boundaries can modify bpe into an incremental patching scheme, but this would increase the length of the byte sequence.}.

\section{{\textsc BLT} Architecture}
\label{section:architecture}
The {\textsc BLT} framework consists of a high-capacity global autoregressive language model, which processes patch-level representations. This global model is complemented by a pair of more compact local models: one encodes byte sequences into patch representations, while the other reconstructs byte sequences from those patch representations~(\autoref{fig:arch_high_level}).

\subsection{Latent Global Transformer Model}

\textit{The Latent Global Transformer} refers to an autoregressive transformer architecture $\mathcal{G}$ composed of $l_{\mathcal{G}}$ layers. This model translates an input series of latent patch representations, $p_j$, into corresponding output patch representations, $o_j$. In this work, the notation $j$ represents indices over patches, while $i$ is reserved for indexing bytes. The global transformer employs a block-causal attention masking scheme~\citep{dubey2024llama}, ensuring that attention at each position can only consider patches up to and including the present patch within a given document. Since the global model is responsible for the majority of computation, both during pre-training and inference, strategically selecting when to activate it enables us to dynamically modulate computational expenditure across different segments of the input and output, depending on their respective complexity.

\subsection{Local Encoder}
\label{section:local}

\textit{The Local Encoder Model}, denoted as $\mathcal{E}$, is implemented as a compact transformer architecture characterized by $l_{\mathcal{E}} << l_{\mathcal{G}}$ layers. Its chief function is to rapidly convert an input sequence of bytes $b_i$ into rich patch-level representations, $p_j$. Notably, this model diverges from the standard transformer design through the integration of a cross-attention module subsequent to each transformer block. This cross-attention operation aggregates byte-level information into patch-level representations, as depicted in~(\autoref{fig:crossattn}).

The transformation begins by mapping the input byte sequence, $b_i$, via an embedding matrix of dimension $\mathbb{R}^{256 \times h_{\mathcal{E}}}$, resulting in the vectors $x_i$. These base embeddings can be further enriched by incorporating optional hash-embeddings, elaborated in~(\S\ref{sec:hash-emb}). The input representations then undergo repeated processing through stacks of interleaved transformer and cross-attention blocks~(\S\ref{sec:enc-cross-attn}), ultimately yielding the patch representations $p_i$ destined for subsequent processing by the global transformer, $\mathcal{G}$.

The attention mechanism within these transformer layers employs a \textit{local block causal} mask, enabling each byte to reference a predetermined window of $w_{\mathcal{E}}$ previous bytes. While these windows may extend across dynamic patch demarcations, they are not permitted to span across distinct document boundaries. Further elaboration on the embedding techniques and the structure of the cross-attention module is provided in the following subsections.

\subsubsection{Encoder Hash n-gram Embeddings}
\label{sec:ngram-emb}
\label{sec:hash-emb}
To ensure that the representations at every step $i$ are both robust and expressive, it is crucial to encode contextual information from the previous bytes. {\textsc BLT} addresses this by representing the current byte $b_i$ both in isolation and within the context of a byte n-gram. Specifically, for each index $i$, we begin by generating byte-grams

\begin{align}
    g_{i,n}=\{b_{i-n + 1},\ldots, b_i\}
\end{align}

for every byte position $i$ and for each $n$ in the range from three to eight.\footnote{
Byte-grams whose length $n$ exceeds $i$ are excluded from consideration. }

Next, we present hash $n$-gram embeddings, which utilize a hash function to associate each possible byte $n$-gram with a position in a corresponding embedding matrix $E_{n}^{hash}$ of predetermined size, for every $n \in \{3,4,5,6,7,8\}$~\citep{bai2010learning}. The embedding derived in this manner is summed with the respective byte’s embedding, subsequently normalized, and then used as the input to the local encoder network. The enhanced embedding is computed as follows:

\begin{align}
e_i &= x_i + \sum_{n = 3,...,8}E_{n}^{hash}(\text{Hash}(g_{i,n})) \\
\text{where, Hash}(g_{i,n}) &= \text{RollPolyHash}(g_{i,n}) \% |E_{n}^{hash}|
\end{align}

We standardize $e_i$ by dividing it by the total number of $n$-gram sizes plus one, and utilize RollPolyHash as specified in Appendix~\ref{appendix:rollpolyhash}. Section~\ref{section:ablations} presents an ablation study examining how varying $n$ in $n$-gram hash embeddings, as well as modifying the size of the embedding table, impact the trends observed in flop-controlled scaling laws. Beyond employing hash-based $n$-gram embeddings, we also explored frequency-based $n$-gram embeddings; the methodology and outcomes of these investigations are described in Appendix~\ref{appendix:ngrams}.

\subsubsection{Encoder Multi-Headed Cross-Attention}
\label{sec:enc-cross-attn}
Our approach largely mirrors the input cross-attention mechanism employed in the Perceiver model~\citep{jaegle2021perceiver}. However, unlike Perceiver, where a static collection of latent vectors is used, our method uses latent representations specific to each variable patch (\autoref{fig:crossattn}), restricting their attention solely to the bytes within their respective patches. In this setup, for each patch $p_j$, a corresponding query vector is constructed by first pooling the representations of all bytes in $p_j$, then mapping this pooled result through a linear layer, $\mathcal{E}_{C} \in \mathbb{R}^{h_{\mathcal{E}} \times (h_{\mathcal{E}} \times U_{\mathcal{E}})}$, where $U_{\mathcal{E}}$ designates the encoder cross-attention head count. Letting $f_{\text{bytes}}(p_j)$ be the sequence of bytes making up patch $p_j$, we then proceed to compute

\begin{align}
P_{0,j} &= \mathcal{E}_{C}(f_\text{bytes}((p_j)), f~ \text{is a pooling function} \\
P_l &= P_{l-1} + W_o\left(\text{softmax}\left(\frac{QK^T}{\sqrt{d_k}}\right)V\right) \\
\text{where } Q_j &= W_q(P_{l-1,j}), K_i = W_k(h_{l-1,i}), V_i = W_v(h_{l-1,i}) \\
h_l &= \text{Encoder-Transformer-Layer}_l(h_{l-1})
\end{align}

Here, $P \in \mathbb{R}^{n_p \times h_{\mathcal{G}}}$ denotes the $n_p$ patch-level representations input to the global model; these representations are initially formed by pooling the byte embeddings $e_i$ associated with each respective patch $p_j$. The matrices $W_q$, $W_k$, $W_v$, and $W_o$ refer to the respective projection matrices for the queries, keys, values, and outputs. Both keys and values are derived by projecting the byte representations $h_i$ output from the preceding layer (or the embeddings $e_i$ for the first layer). A customized masking approach is applied so that each query $Q_j$ is restricted to attending only to keys and values associated with bytes from the same patch $j$. Due to the use of multi-head attention across $Q$, $K$, and $V$, and because patch representations typically have a higher dimensionality ($h_{\mathcal{G}}$) than $h_{\mathcal{E}}$, we keep $P_l$ structured as multiple heads of $h_{\mathcal{E}}$-dimensional representations during cross-attention, concatenating these later to produce $h_{\mathcal{G}}$-dimensional outputs. Further, we adopt pre-LayerNorm for the queries, keys, and values, and the cross-attention mechanism does not employ positional embeddings. The cross-attention block also incorporates a residual connection.

\subsection{Local Decoder} 

The local decoder $\mathcal{D}$, akin to the local encoder in its design, is implemented as a lightweight transformer model comprising $l_{\mathcal{D}}$ layers, with $l_{\mathcal{D}}$ being significantly fewer than $l_{\mathcal{G}}$. Its role is to convert a sequence of global patch representations, $o_j$, back into raw bytes, denoted as $y_i$. Operating autoregressively, the local decoder generates each byte based on the sequence of previously produced bytes, leveraging the hidden representations generated by the local encoder as inputs for the byte sequence. The architecture employs $l_{\mathcal{D}}$ layers in an alternating pattern between cross-attention modules and transformer blocks. Within each decoder layer, cross-attention is applied before the transformer operation, allowing the model to construct representations for each byte from the input patch representations, after which the transformer layer processes the output byte sequence.

\subsubsection{Decoder Multi-headed Cross-Attention} 

In the decoder cross-attention, the roles of the queries and key/values are interchanged i.e. the byte-representations are now the queries, and the patch representations are now the key/values. The initial byte-representations for the cross-attention are initialized as the byte embeddings from the last encoder layer i.e. $h_{l_{\mathcal{E}}}$. The subsequent byte-representations for layer $l$, $d_{l,i}$ are computed as:

\begin{align}
D_0 &= h_{l_{\mathcal{E}}} \\
B_l &= D_{l-1} + W_o\left(\text{softmax}\left(\frac{QK^T}{\sqrt{d_k}}\right)V\right), \\
  &\text{where } Q_i = W_q(d_{l-1,i}), K_i = W_k(\mathcal{D}_C(o_j)), V_i = W_v(\mathcal{D}_C(o_j)) \\
  D_l &= \text{Decoder-Transformer-layer}_l(B_l)
\end{align}

In this context, $W_k$ and $W_v$ represent key and value projection matrices responsible for performing a linear transformation followed by the split operation $\mathcal{D}_C$ on the final patch representations $o_j$ generated by the global model. The query projection matrix $W_q$ acts on byte representations $d_{l-1}$ produced by the preceding decoder transformer layer, or $h_{l_{\mathcal{E}}}$ in the case of the initial layer. Meanwhile, $W_o$ is defined as the output projection matrix, which ensures $B$ has dimensions $\mathbb{R}^{h_{\mathcal{D}} \times n_b}$, where $n_b$ denotes the number of produced bytes. The subsequent decoder representations $D_l$ are then produced by applying a decoder transformer layer to the result of the cross-attention block, $B$. Following the approach adopted in the local encoder’s cross-attention mechanism, our architecture incorporates multiple attention heads, employs pre LayerNorms, omits positional embeddings, and utilizes a residual connection encasing the cross-attention module.

\section{Experimental Setup}
\label{section:experiments}
We established a suite of controlled experiments designed to rigorously evaluate {\textsc BLT} against models employing tokenization, explicitly ensuring that {\textsc BLT} does not receive any undue benefits from the potential use of extended sequence lengths.

\subsection{Pre-training Datasets} In all our model size explorations, we employ two primary pre-training datasets: (1) The Llama 2 dataset~\citep{touvron2023llama}, consisting of 2 trillion tokens sourced from a diverse range of public data, which has undergone rigorous cleaning and filtering to enhance data quality; and (2) {\textsc BLT}-1T, a newly assembled corpus totaling 1 trillion tokens compiled from assorted publicly accessible sources and incorporating a subset of the pre-training set made available through Datacomp-LM~\citep{li2024datacomp}. The Llama 2 dataset is leveraged for scaling law analyses—specifically, those pertaining to optimal token count as outlined by~\cite{dubey2024llama}—to guide architectural design for {\textsc BLT}. In contrast, {\textsc BLT}-1T is utilized in a full pre-training cycle aimed at evaluating downstream task performance relative to Llama 3. Importantly, neither dataset includes data originating from Meta products or services. In addition, our tokenizer-based model baselines employ the Llama 3 tokenizer (with a 128K-token vocabulary), as it consistently delivered improved baseline results compared to the Llama 2 tokenizer within our empirical studies.

\subsection{Entropy Model} For our experiments, the entropy model is implemented as a byte-level language model, which is trained using the same data distribution as the main {\textsc BLT} model. By default, we utilize a transformer architecture consisting of 100M parameters, 14 layers, and a hidden dimension of 512, employing a sliding window attention span of 512 bytes. All other hyperparameter configurations are matched to those used in our local and global transformers. We have explored a range of model capacities, context window sizes, and architectural variations as outlined in \autoref{section:ablations}. Notably, if the model's receptive field is sufficiently restricted, the resulting entropy model is compact enough to be represented as a high-efficiency lookup table.

\subsection{Entropy Threshold and Equalizing Context Length} For models utilizing entropy-driven patching, we determine a patching threshold so as to attain a specified average \textit{patch size} when evaluated on the pretraining data mixture. In {\textsc BLT}, unlike traditional tokenization approaches, the choice of \textit{patch size} is flexible and can be set arbitrarily, which substantially influences the effective context length for the model. To guarantee that different patch sizes do not result in discrepancies in accessible context—thereby avoiding undue advantages—we fix the expected number of bytes per batch. Consequently, for configurations employing larger patch sizes, the number of patches (i.e., the sequence length) per batch is proportionally decreased. For models trained on Llama 2 data, we adopt an 8k byte context. In contrast, for training on the {\textsc BLT}-1T dataset, the average context window is raised to 16k bytes, all while keeping the average batch size at 16M bytes consistent across configurations.

Although the mean batch size does not change, dynamic patching approaches generate varying proportions of bytes per patch when forming data batches. To enhance computational efficiency, our implementation of {\textsc BLT} training assembles batches by grouping patches directly, thereby bypassing the need for additional padding in the resource-intensive latent transformer. This batching strategy guarantees a fixed patch count for each batch. To further manage computational load, we pad or, if necessary, truncate the input byte sequences to 12k bytes for Llama 2 and 24k bytes for the {\textsc BLT}-1T datasets during training. This practice prevents excessive memory consumption that could occur due to exceptionally large patch sizes within certain sequences.

\subsection{Entropy Model Context}
\label{section:entropy-context}
Our observations indicate that applying entropy patching leads to increasingly larger patches, particularly in highly structured formats such as multiple choice tasks, which tend to exhibit significant repetition (as illustrated by the patching process for an MMLU example in \autoref{fig:mmlu_patching}). This outcome arises because the repeated segments in the entropy model context exhibit lower entropy. To address this in the large-scale execution of {\textsc BLT}-Entropy with a patch size of 4.5, we refresh the entropy context by inserting new lines and employ an approximate monotonicity constraint, since this approach mitigates "entropy drift" due to variations in context length. Notably, this adjustment solely alters our entropy calculation methodology, while our procedure for selecting the entropy threshold remains unchanged.

\subsection{FLOPs Estimation} 
\label{section:flops}

Our methodology for estimating transformer flops primarily adheres to the formulations established in Chinchilla~\citep{hoffmann2022training}, which include the computational costs associated with feed-forward networks, qkvo projections within the self-attention module, attention matrix calculations, and the final output projection. A significant deviation from prior conventions is our assumption that the input embedding layer utilizes an optimized lookup mechanism rather than conventional dense matrix multiplication, effectively rendering it a zero-flop operation. Consistent with prior research, we also account for the backward pass by assuming it incurs twice the flop count of the forward pass.

To compute flops \textit{per byte} for {\textsc BLT} models, we add up the flops for the local encoder transformer, the global latent transformer, and the local decoder transformer, together with the cross attention blocks in the encoder and the decoder:

\begin{align}
    \text{FL}_{\text{{\textsc BLT}}} &= \frac{1}{n_p}\,\text{Transf. FL}(h_{\mathcal{G}}, l_{\mathcal{G}}, m=n_{ctx}/n_p, V=0)\\
   &\quad +\,\text{Transf. FL}(h_{\mathcal{E}}, l_{\mathcal{E}}, m=w_{\mathcal{E}}, V=0) \\
   &\quad +\,\text{Transf. FL}(h_{\mathcal{D}}, l_{\mathcal{D}}, m=w_{\mathcal{D}}, V=256) \\ 
    &\quad +\, k/n_p \cdot \text{Cross Attn. FL}(h_{\mathcal{E}}, l_{\mathcal{E}}, m=n_p, r=n_p/k) \\ 
   &\quad +\, \text{Cross Attn. FL}(h_{\mathcal{D}}, l_{\mathcal{D}}, m=k, r=k/n_p) 
\end{align}

Here, $n_{ctx}$ denotes the length of the sequence in bytes, $n_p$ specifies the size of each patch, $r$ represents the proportion of queries to keys/values, and $k$ indicates the ratio between patch dimension and byte dimension; that is, $k$ reflects how many segments of the local model are merged to create the full model representation (with $k=2$ in \autoref{fig:crossattn}). The variable $V$ stands for the output projection’s vocabulary size and is utilized exclusively in the local decoder. The effective context length used by the attention mechanism, denoted as $m$, varies based on whether the operation acts on the sequence of bytes or patches. We adjust the attention flop calculations separately for each part accordingly. Detailed equations outlining how Transformer-FLOPs and Cross-Attention FLOPs are calculated appear in Appendix~\ref{section:flops-appendix}.

\subsection{Bits-Per-Byte Estimation} Perplexity only makes sense in the context of a fixed tokenizer as it is a measure of the uncertainty for each token. When comparing byte and token-level models, following previous work~\citep{xue2022byt5, yu2023megabyte, wang2024mambabyte}, we instead report Bits-Per-Byte (BPB), a tokenizer independent version of perplexity. Specifically:

\begin{align}
    \text{BPB}(x)&= \frac{\mathcal{L}_{CE}(\pmb{x})}{\ln(2)\cdot n_{\text{bytes}}}
\end{align}

In this expression, the aggregate cross-entropy loss that quantifies the uncertainty associated with the data $\pmb{x}$ is divided by the product of a normalization constant and the total number of bytes in $\pmb{x}$.

\subsection{Transformer Architecture Hyperparameters} In configuring all transformer modules within {\textsc BLT}—encompassing both its local and global components—we adhere primarily to the specifications established in the Llama~3~\citep{dubey2024llama} framework. Specifically, feed-forward networks utilize the SwiGLU activation function~\citep{swiglu}, while self-attention mechanisms employ rotary positional embeddings (RoPE)~\citep{RoPe} with parameter $\theta=500000$~\citep{xiong2024effective}. Layer normalization across the model is carried out using RMSNorm~\citep{RMSNorm}. For every self-attention layer with conventional attention patterns such as \textit{block causal} or \textit{fixed-window block causal}, Flash attention~\citep{dao2022flashattention} is implemented, setting the window size at 512 for fixed-width attention configurations. Owing to the dynamically changing, patch-specific masks utilized in cross-attention layers, we employ Flex Attention\footnote{\url{https://pytorch.org/blog/flexattention}}, enabling fused computation and substantial improvements in training efficiency.

\subsection{{\textsc BLT}-Specific Hyperparameters} To assess the performance of {\textsc BLT} models, our experimental setup investigates both scaling behavior and evaluation on downstream tasks, encompassing a range of model sizes: 400M, 1B, 2B, 4B, and 8B. Appendix~Table~\ref{tab:arch} details the specific architecture hyperparameters used across these configurations. 

For the initial cross-attention layer within the local encoder, query vectors are initialized using max-pooling. Every {\textsc BLT} model variant is configured with $500,000$ hashes, employing a single hash function, and utilizes n-gram sizes from 3 up to 8. A uniform learning rate of $4e-4$ is adopted throughout all model scales. This setting aligns the learning rate for both token-based and {\textsc BLT} architectures, as a targeted hyperparameter search (conducted on the 400M and 1B models) spanning $1e-3$ to $1e-4$ indicated no benefit in adopting distinct rates. When scaling on Llama-2 datasets, we adhere to the batch-size guidelines from~\cite{dubey2024llama} or use an equivalent amount measured in bytes. For the optimization procedure, we rely on the AdamW optimizer~\citep{adamw}, configuring it with $\beta_1=0.9$, $\beta_2=0.95$, and setting $\epsilon=10^{-8}$. A linear warm-up is applied over the initial 2000 steps, after which the learning rate decays to zero following a cosine schedule. We also enforce a weight decay of 0.1 and restrict the global gradient norm to a maximum of 1.0 through gradient clipping.

\section{Scaling Trends}
\label{section:scaling}

We offer a comprehensive overview of the scaling behaviors observed in byte-level models, providing guidance for advancing the scalability of {\textsc BLT} models. This scaling analysis is designed to overcome several shortcomings found in previous explorations of byte-level architectures: (a) We examine trends within a compute-optimal training framework, (b) We develop corresponding 8B-parameter models using substantial datasets, reaching up to 1T tokens/4T bytes, and assess their performance across downstream benchmarks, and (c) We evaluate scaling dynamics under fixed inference-cost scenarios. In a subsequent section, we will delve into the particular strengths afforded by modeling at the byte-sequence level.

\subsection{Parameter Matched Compute Optimal Scaling Trends}
\label{section:parameter-matched-scaling}
We utilize the Llama 2 dataset to train both \textit{compute-optimal} bpe and {\textsc BLT} models across four distinct parameter regimes, specifically from 1B up to 8B parameters. Training flops are plotted against language modeling performance, evaluated on a representative subset of the data mixture used during training. The bpe models follow the model-to-data ratio identified as optimal in Llama~3~\citep{dubey2024llama}. This approach is considered \textit{compute-optimal}, theoretically maximizing performance on the training corpus under a specified computational constraint~\citep{hoffmann2022training}, and serves as a strong baseline in our comparisons. Correspondingly, for every bpe model, an equivalent {\textsc BLT} model is trained on identical data, employing a Latent Transformer architecture that mirrors the size and structural configuration of its bpe Transformer counterpart.

As shown in~\autoref{fig:scaling-compute-opt} (right), {\textsc BLT} models consistently perform on par with, or surpass, their bpe-based equivalents, and this pattern remains consistent as both model size and compute (FLOPs) increase. To our knowledge, {\textsc BLT} represents the first byte-level Transformer architecture to demonstrate scaling behavior similar to BPE-based models within compute-optimal settings. This observation supports our hypothesis that the optimal parameter-to-compute ratio found for bpe architectures is also applicable to {\textsc BLT}, or at minimum, does not significantly deviate from it.

Enhancements to model architectures, alongside the adoption of dynamic patching, are both essential to effectively align with bpe scaling patterns. In ~\autoref{fig:scaling-compute-opt} (left), we present a comparison between models utilizing space-patching approaches and Llama 3. To estimate the performance of SpaceByte \citep{slagle2024spacebyte}, we employ {\textsc BLT} space-patching configured without n-gram embeddings or cross-attention. While SpaceByte represents a substantial advancement over Megabyte, its performance still lags considerably behind Llama 3. \autoref{fig:scaling-compute-opt} (right) demonstrates the collective gains achieved through refined architectural designs as well as dynamic patching methods. At scale, {\textsc BLT} architectures are competitive with leading tokenizer-based models, such as Llama 3.

Additionally, our findings indicate that, for tokenizer-based models, the selection of tokenizer significantly impacts performance; specifically, employing the Llama-3 tokenizer leads to superior results compared to utilizing the Llama-2 tokenizer, even when the training dataset remains unchanged.

In summary, our {\textsc BLT} architecture exhibits behavior intermediate to that of Llama 2 and 3, particularly when much larger patch sizes are employed. While the bpe tokenizers for Llama 2 and 3 possess average token lengths of 3.7 and 4.4 bytes, respectively, {\textsc BLT} demonstrates analogous scaling properties even at average patch sizes of 6 or 8 bytes. Since inference flops are inversely related to the mean patch size, adopting a patch size of 8 bytes could yield close to a 50\% reduction in inference computational cost. Moreover, increasing patch sizes appears to enhance model performance as both model and dataset sizes are scaled up. Notably, {\textsc BLT} using an 8-byte patch size initially underperforms relative to the bpe-based Llama 2 at the 1B parameter mark, but surpasses it once scaled to 7B parameters. This implies that employing such large patch sizes could result in further performance gains at even greater scales, and that utilizing even larger patches may be attainable as model parameters and available training compute increase.

\subsection{Beyond Compute Optimal Task Evaluations}
\label{section:evals}
To explore the scaling behavior in greater depth, we train an 8B {\textsc BLT} model past the compute-optimal threshold using the {\textsc BLT}-1T dataset, which is both larger in size and of higher quality. We then evaluate its performance across a diverse set of widely-used benchmarks for classification and generation tasks. Our evaluation spans benchmarks covering common sense reasoning, world knowledge, as well as code generation capabilities:
\paragraph{Classification tasks} we consider ARC-Easy (0-shot)~\citep{Eval_ARC}, Arc-Challenge (0-shot)~\citep{Eval_ARC}, HellaSwag (0-shot)~\citep{Eval_hellaswag}, PIQA (0-shot)~\citep{Eval_piqa}, and MMLU (5-shot)~\citep{hendrycks2020measuring}. We utilize a prompt-scoring approach—computing the probability assigned to each candidate choice—and present mean accuracy across these datasets.
\paragraph{Coding related generation tasks:} pass@1 metrics are reported for MBPP (3-shot)~\citep{Eval_MBPP} and HumanEval (0-shot)~\citep{Eval_HumanEval} to assess the proficiency of the language models in Python code synthesis.

In \autoref{tab:evals}, we present a comparison among three models trained using the {\textsc BLT}-1T dataset: a model leveraging a bpe Llama 3 tokenizer,\footnote{The Llama 3 tokenizer, which provides a 128k vocabulary, is preferred due to its superior performance relative to the 32k-vocabulary Llama 2 tokenizer.} and two variations of the {\textsc BLT} model—one that adopts a space-patching approach ({\textsc BLT}-Space) and another that employs an entropy-driven patching technique ({\textsc BLT}-Entropy). The entropy-based variant is further constrained to approximate monotonicity and resets its context on encountering new lines, as outlined in~\autoref{section:entropy-context}. Each model is trained under comparable flop budgets. Notably, for {\textsc BLT}-Entropy, we introduce an inference phase modification by lowering the entropy threshold from 0.6 to 0.1. This adjustment enhances performance across evaluated tasks, although it increases the number of inference steps required.

The {\textsc BLT}-Entropy model achieves superior results compared to the Llama 3 model on 4 of the 7 evaluated tasks, despite both models being trained with an equivalent quantity of bytes. This performance gain likely stems from two main factors: (1) more efficient utilization of training computation through dynamic patching, and (2) the explicit representation of information at the byte level rather than at the token level.

Conversely, {\textsc BLT}-Space delivers weaker results than the Llama 3 tokenizer across all tasks except one. Nevertheless, due to its greater mean patch size of 6 bytes, it notably decreases inference flops. By contrast, both the bpe and entropy-patching models display similar mean patch sizes, close to 4.5 bytes on the mixed training dataset. When subjected to an identical training budget, the model with the larger patch size is able to process 30\% more data compared to the other two, potentially moving {\textsc BLT} further from the point of computational optimality.

\subsection{Patches Scale Better Than Tokens} Using {\textsc BLT} models, it is possible to enlarge both the model and the patch size simultaneously without increasing either the training or inference computational costs, provided that the training dataset remains unchanged. This scalability is distinctive to patch-based architectures, which, unlike fixed-vocabulary token-based models, allow the patch size to grow arbitrarily without the same efficiency limitations (see Section~\ref{section:bpe-patch}). Employing longer patches reduces the computational burden, enabling those resources to be used for expanding the global latent transformer since it processes input less frequently.

In this work, we perform a fixed inference scaling analysis to evaluate whether increasing model size while decreasing the number of inference steps on larger input patches yields better performance compared to smaller models operating with more steps. Utilizing initial model configurations of 400m and 3.6B parameters and the Llama 2 tokenizer, we identify flop-matched models employing the Llama 3 tokenizer and {\textsc BLT}-Entropy models featuring mean patch sizes of 6 and 8 bytes when trained on the datamix dataset (refer to~\autoref{tab:fixed-inf-params} for specific model configurations). For those models using patch size 8, we increase the number of encoder layers to 3 rather than 1. Each model is trained across a range of training flop budgets.

\autoref{fig:fixed_inference_scaling} illustrates that {\textsc BLT} models exhibit superior scaling behavior compared to tokenization-based models across both inference flop categories. For both scenarios, BPE-based architectures outperform at low training compute levels, but {\textsc BLT} models soon overtake them once training scales past the compute-optimal threshold. This suggests that allocating additional resources during initial pretraining can yield higher-quality models while maintaining a constant inference compute limit. A notable case is the 8B parameter category, exemplified by Llama 3.1, which underwent pretraining on datasets two orders of magnitude greater than what would be deemed compute-optimal for its size.

The point at which {\textsc BLT} outperforms token-based models now aligns more closely with the compute-optimal regime as model size increases (shifting from 3x to 2.5x the compute-optimal budget for the larger flop class models). In the same vein, the patch size 8 model in the larger flop category demonstrates a more pronounced scaling trajectory, overtaking other model variants at an earlier stage. As outlined in Section~\ref{section:parameter-matched-scaling}, scaling models with larger patch sizes tends to narrow the gap in performance compared to BPE models at higher capacities. We partially explain this effect by the reduction in the proportion of total flops allocated to the byte-level Encoder and Decoder components, whose computational growth lags behind that of the Latent Transformer. When scaling up the model from 400M to 8B parameters (a 20x increase), the local parameters specific to {\textsc BLT} increase by only about two times. Notably, adjusting the patch size affects the flops contributed by the patch Latent Transformer without impacting those from the byte-level components. This dynamic explains why the model size ratio for {\textsc BLT}-Entropy with ps=8 increased slightly, from 1.6x to 1.7x relative to the Llama 2 baseline, when moving to higher model scales.

In conclusion, our investigation into scaling patch lengths reveals that the {\textsc BLT} patch-based architecture exhibits superior scaling behavior when both patch size and overall model size are increased in tandem. These positive trends appear to hold and further improve as the model becomes larger.

\section{Byte Modeling Improves Robustness} We further evaluate how the robustness of {\textsc BLT} compares with that of token-based models that do not encode direct byte-level information, and introduce a method for converting pretrained token-based models to operate on bytes. 
\label{sec:robustness}

\subsection{Character-Level Tasks}

Initial efforts to develop byte-level models were largely driven by the desire to leverage their inherent resilience to noise at the byte level in input data, as well as their capability to comprehend the internal structure of tokens—areas where traditional tokenizer-based approaches often fall short. To assess these attributes, we conduct further experiments using benchmarks specifically designed to gauge both resistance to noisy input and sensitivity to character-level features, encompassing English and multiple languages, as well as incorporating digits and phonemes. The outcomes of these tests are summarized in Table~\ref{tab:char_tasks}.

\paragraph{Noisy Data} To assess and compare the robustness of tokenizer-based approaches versus {\textsc BLT}, we construct corrupted versions of the benchmark classification datasets outlined in Section~\ref{section:evals}. Specifically, we apply five varied character-level noising methods to the text: (a) \textit{AntSpeak}: All text is transformed into uppercase and each character is separated by a space. (b) \textit{Drop}: Eliminates 10\% of characters from the string at random positions. (c) \textit{RandomCase}: Randomly changes the case of each character so that approximately half are uppercase and half are lowercase. (d) \textit{Repeat}: Increases the frequency of 20\% of characters, repeating them up to four times. (e) \textit{UpperCase}: Converts the complete text to uppercase letters. Each noise procedure is individually applied during testing to either the prompt, the response, or both—considered as independent tasks—and we record the mean accuracy over these configurations. Table \ref{tab:char_tasks} provides performance outcomes on noised HellaSwag~\citep{Eval_hellaswag}; these results show that {\textsc BLT} consistently exhibits greater resilience than tokenizer-based systems, surpassing them by an average of 8 points, and even exceeds the performance of the Llama 3.1 model despite the latter being trained on significantly more data.

\paragraph{Phonology - Grapheme-to-Phoneme (G2P)} We evaluate the proficiency of {\textsc BLT} in converting a sequence of graphemes—that is, the individual characters forming a word—into its corresponding phonemic transcription, effectively capturing the word’s spoken form. Table \ref{tab:char_tasks} details the outcomes of the G2P evaluation, which was conducted under a 5-shot experimental regime leveraging Phonology Bench~\citep{suvarna2024phonologybench}. The results show that {\textsc BLT} consistently surpasses the performance of the baseline model based on the Llama 3 1T tokenizer for this specific task.

\paragraph{CUTE}
To evaluate character-level comprehension, we test {\textsc BLT} on the CUTE benchmark~\citep{edman2024cute}, which includes a range of tasks grouped into three main areas: compositionality, orthographic similarity, and sequence manipulation skills. The benchmark is particularly challenging for most models relying on tokenization, which seem to encode knowledge of token spellings but have difficulty applying this understanding for direct textual manipulation. As presented in Table~\ref{tab:char_tasks}, {\textsc BLT}-Entropy surpasses both BPE Llama 3 models by a margin of over 25 points on this suite of tasks. Notably, our model excels in tasks requiring character-level manipulations, reaching 99.9\% accuracy on spelling assessments. The magnitude of improvement is remarkable given that {\textsc BLT} was trained with only a sixteenth of the data used for Llama 3.1, underscoring the inherent challenges BPE-based models face in learning character-level representations. Figure \ref{fig:cute_examples} provides examples where Llama 3 with BPE tokenization falls short while {\textsc BLT} delivers strong results. BPE models only show an advantage in word insertion and deletion; these specific word-level edits can be less trivial for byte-level models, yet the performance gap remains modest, and learning to compose words from characters could prove easier than the reverse. We adopt the identical evaluation strategy and original Huggingface prompts for all models. Additional prompt design could potentially enhance the performance of BPE-based architectures.

\paragraph{Low Resource Machine Translation} We assess {\textsc BLT} on both source and target sides for six widely-spoken language families and twenty-one lesser-resourced languages in diverse scripts, following the FLORES-101 benchmark~\citep{goyal2022flores}. SentencePiece BLEU scores are shown in Table~\ref{tab:flores}. Our findings reveal that {\textsc BLT} surpasses a Llama 3 tokenizer-based system, gaining a 2-point margin for translations into English and a 0.5-point improvement in translations from English. Across high-resource language pairs, {\textsc BLT} matches or modestly outperforms Llama 3. Notably, in many cases within lower-resource language families, {\textsc BLT} yields superior results, highlighting the strengths of byte-level modeling when generalizing over long-tail byte patterns.

\subsection{Training {\textsc BLT} from Llama 3}
\label{sec:distillation}
We investigate an approach in which {\textsc BLT} models can capitalize on pre-existing, pretrained tokenizer-based models to enhance and accelerate training convergence. This is accomplished by initializing the global transformer parameters of {\textsc BLT} with those from a pretrained Llama 3.1 model. The training procedure then updates the global transformer weights at a rate one-tenth that of the local encoder and local decoder learning rates, maintaining the same optimal step count used for Llama 3. We summarize results against a standard {\textsc BLT} configuration in Table~\ref{tab:distill}. The findings indicate that the {\textsc BLT} initialized with Llama 3.1 substantially surpasses both the standard Llama 3 and {\textsc BLT} baselines, when training resources (flops) are held constant. Notably, this initialization strategy allows the model to outperform our {\textsc BLT}-Entropy variant (refer to Table~\ref{tab:evals}), which was exposed to a much more extensive training corpus (1T tokens), on the MMLU benchmark. These observations suggest that initializing from Llama 3.1 is a highly effective method for considerably reducing computational cost during training.

This approach essentially repurposes tokenizer-based architectures as tokenizer-free models, thereby transforming a pre-trained LLaMA 3.1 model into a {\textsc BLT} configuration. For a thorough comparison, we report results for the original LLaMA 3.1—trained on 15T tokens—in Table \ref{tab:distill}, contrasting its performance with that of the LLaMA-based {\textsc BLT}. Although our model incurs only marginal losses on MMLU and HumanEval, it exhibits more pronounced declines in performance on several other benchmarks. These findings indicate that additional research is required to make optimal use of the pre-trained model, specifically through refining data mixture strategies and tuning other key hyperparameters.

\section{Ablations and Discussion}
\label{section:ablations}
Here, we present a series of ablation studies that provide rationale for the architectural decisions underlying {\textsc BLT} as well as the patching strategy and selection of hyper-parameters employed in the {\textsc BLT} model with 8B parameters, which was trained using the {\textsc BLT}-1T dataset.

\paragraph{Entropy Model Hyper-parameters} To investigate how the capacity of the entropy model and the size of the context window impact scaling behavior, we conduct experiments by training byte-level entropy transformer architectures across a range of sizes, spanning from 1 million to 100 million parameters, and employing various context window lengths between 64 and 512 tokens. We illustrate the scaling laws by plotting bits-per-byte (bpb) against training FLOPs for our $400m$ and $1b$ {\textsc BLT} models, each trained on the Llama-2 dataset, as shown in \autoref{fig:entropy_ablation}. Our findings suggest a clear positive association between scaling outcomes and both the model size and the context window length, although the advantages become marginal when expanding beyond 50 million parameters.

\paragraph{Types of Patching}

We conduct an ablation study on the four patching strategies outlined in Section~\ref{section:patching}: (1) Strided Patching using strides of 4 and 6, (2) Whitespace-based Patching, (3) Patching through the BPE Tokenizer associated with the Llama 3 tokenizer, and (4) Entropy-driven patching guided by a compact byte-level LLM.

Although dynamic patching shortens the sequences' actual length, we ensure that the context length remains consistent across all patching approaches by explicitly controlling for sequence length. Consequently, during both training and inference, each model is exposed to the same expected number of bytes per sequence, thereby eliminating potential confounds related to the capacity for modeling longer contexts. Figure~\ref{fig:scaling-compute-opt} presents the outcomes of these ablation studies. Notably, every alternative patching method surpasses static patching, with space patching performing nearly as well as dynamic entropy based patching.

\autoref{tab:evals_ablation} summarizes benchmark results for various {\textsc BLT} models, offering a comparison among tokenizer-based approaches, space patching, and entropy-based patching, all of which were trained on the Llama 2 dataset for a tuned number of steps~\citep{dubey2024llama}. While space patching serves as a more straightforward method—eschewing the need to invoke an entropy model during training—our findings demonstrate that the improvements attributed to entropy-based patching, previously identified in the context of scaling behavior~(Section~\ref{section:scaling}), also persist across subsequent benchmark tasks.\footnote{Results for space patching originate from prior experiments conducted without cross-attention, yet comparable trends are observed with the inclusion of cross-attention.}

\paragraph{Cross-Attention} 
In~\autoref{tab:crossattn_ablations_1b}, we investigate the effect of adding cross-attention mechanisms at different stages within both the encoder and decoder components of {\textsc BLT}. For encoder cross-attention, we evaluate three approaches for initializing the queries: 1) utilizing an identical trainable embedding for all global states, 2) applying a hash embedding derived from the patch's bytes, and 3) employing a pooling operation over the encoder’s hidden representations for the patch bytes at the specific encoder layer.

Our results indicate that incorporating cross-attention within the \textit{decoder} yields the highest effectiveness. In contrast, employing cross-attention in the encoder leads to marginal gains, and these occur solely when queries are initialized via pooling. Moreover, our analysis shows that the benefits of cross-attention are particularly evident on Common-Crawl data and become more pronounced as patch sizes increase.

\paragraph{n-gram Hash Embeddings} 

We experiment with n-gram hash embedding vocabularies of sizes 0, 100K, 200K, and 400K, and display the corresponding outcomes in Table~\ref{tab:ngram_ablations_1b}. Our findings suggest that hash embeddings contribute positively across all domains, with the most notable impact observed on Wikipedia and Github datasets (yielding a 0.04 bpb difference, versus 0.01 bpb after 15k steps at the 8B scale). When scaling to 8B parameters, reducing the hash count from 500K to 300K leads to only a minor change in performance—approximately 0.001 bpb after 15k steps. This result suggests that hash embeddings play a critical role in enabling {\textsc BLT} models to achieve competitive performance with those using traditional tokenizers, although performance gains plateau beyond 300K hashes. Furthermore, the improvements appear to act in a complementary fashion with cross-attention, as benefits are observed across distinct datasets.

\paragraph{Local Model Hyperparameters} Table \ref{tab:local_layers} presents an ablation study on different configurations for the number of layers within the local encoder and decoder. Results indicate that, particularly when utilizing hash n-gram embeddings, {\textsc BLT} achieves strong performance with a minimally complex encoder—specifically, a single layer—while performance benefits from employing a more substantial decoder.

\section{Related Work}
\label{section:related_work}

\textbf{Character-Level RNNs:} The task of Character Language Modeling has maintained popularity since the advent of neural models~\citep{sutskever2011generating,mikolov2012subword,graves2013generating}, primarily due to its inherent capacity to naturally represent words not found in the training vocabulary, without depending on back-off mechanisms. In related work, \cite{kim2016character} developed a model that encodes input exclusively at the character level using a combination of convolutional and highway networks, followed by LSTM-based RNNs. Their model achieved parity with the then state-of-the-art RNN language models in English and surpassed them for morphologically complex languages—highlighting another frequently cited benefit of character-based LLMs. Similarly, \cite{kenter2018byte} employed byte-level LSTM architectures for machine comprehension tasks and observed superior results compared to word-level approaches, particularly for languages with intricate morphology such as Turkish and Russian. In a comparable direction, \cite{zhang2015character} introduced character-level convolutional models for classification, which outperformed word-based counterparts on some benchmarks. The approach proposed by \citet{chung2019hierarchical} utilizes hierarchical LSTM architectures augmented with boundary detectors at each layer, enabling the model to uncover the latent hierarchical structure of text and thus advance character-level language modeling performance. Diverging from attention mechanisms, ByteNet by \cite{kalchbrenner2016neural} applies convolutional neural network layers on character sequences for the purpose of machine translation.

\textbf{Character-Level Transformers:} Transformer architectures that leverage attention mechanisms~\citep{vaswani2017attention} in combination with subword tokenization~\citep{sennrich-etal-2016-neural} have led to substantial gains in neural language modeling and on standard benchmarks. Nevertheless, employing word and sub-word tokens imposes an inherent inductive bias regarding the model's abstraction level. In an effort to bridge successful transformer methodologies with the positive indications from character-level language models, \citet{al2019character} implemented deep transformer networks augmented by auxiliary loss functions. This enabled their transformer-based approaches to surpass the performance of former LSTM-based character language models. Despite this progress, a notable performance disparity remained when compared to large language models trained on word-level representations. Additionally, GPT-2~\citep{radfordgpt2} reported that, for extensive datasets such as the 1 billion word benchmark, byte-level language models failed to match the effectiveness of their word-level counterparts.

Although \cite{Choe2019BridgingTG} showed that transformer-based LLMs processing data at the byte level can surpass the performance of similarly sized subword-level LLMs, these models require significantly more computational resources and incur longer training times. In a parallel vein, \cite{el2020characterbert} introduced CharFormer—a variant of BERT employing convolutions on character-level embeddings to form word representations—reporting performance gains within the medical domain, albeit at the expense of greater computational costs. CANINE, developed by \cite{clark2022canine}, is a 150M parameter encoder-only architecture that ingests raw character sequences. The core of CANINE consists of a deep transformer, analogous to our proposed global model, supplemented by a local transformer and strided convolutional layers for input downsampling. On various multilingual downstream tasks, CANINE exceeds the performance of mBERT, a corresponding token-level encoder-only model. In ByT5~\citep{xue2022byt5}, the authors investigate byte-level encoder-decoder designs eschewing any form of patching; their approach improves noise robustness and achieves competitive results relative to tokenizer-based models using only a quarter of the training data. Nevertheless, omitting patching operations obligates the model to perform costly attention calculations over each byte, dramatically increasing the computational burden. Modeling input at the byte level rather than with subwords intrinsically extends sequence lengths, presenting significant challenges to scaling such architectures efficiently. More recently, leveraging the Mamba Architecture~\citep{gu2023mamba}—which sustains fixed-size memory states even across extended context windows—\cite{wang2024mambabyte} implemented a patchless, byte-level Mamba model. Their results demonstrate that, at 350M parameters and under fixed FLOP budgets, this Mamba-based model can outperform comparable byte-level transformers on bits-per-byte metrics across multiple datasets.

\textbf{Patching-based approaches:} Leveraging patching strategies can reduce the high computational demands typically associated with byte-level LLMs, while still preserving their performance. Numerous studies have shown promising outcomes in limited settings involving smaller models and training datasets. For example, \cite{nawrot-etal-2022-hierarchical} explore the application of static patching for downsampling and upsampling, introducing the hourglass transformer architecture, which surpasses other byte-level models at the 150M scale. Subsequent work by \cite{nawrot-etal-2023-efficient} advances this further through dynamic patching mechanisms. They present a variety of techniques such as an end-to-end trained boundary-predictor, a predictor supervised with specific tokenizers, and an entropy-driven patching model related to {\textsc BLT}. Their results indicate that these dynamic approaches can exceed the performance of standard transformer models in language modeling benchmarks, particularly at a 40M parameter size on approximately 400M tokens. In addition, \cite{lester2024training} analyze the effects of training on input sequences that have been compressed using arithmetic coding, achieving a level of compression superior to BPE. By implementing an equal-info windows strategy, they achieve better performance than byte-level baselines for language modeling, yet still fall short of those based on subword units.

Our research is primarily motivated by, and shares close ties to, MegaByte~\citep{yu2023megabyte}, a decoder-only causal LLM architecture which employs a predetermined patching strategy and concatenates representations in order to map bytes into fixed-sized patches; on the decoder side, it incorporates a local module to transform patches back to the original byte sequence. The authors of MegaByte show that their approach achieves performance comparable to tokenizer-based models at the 1B parameter scale when evaluated on a dataset comprising roughly 400B bytes. We incorporate a thorough ablation of MegaByte in each of our experimental evaluations, noting that static patching falls short of surpassing state-of-the-art, compute-optimally trained models relying on tokenizers, particularly when computation is controlled. In contrast, we illustrate how {\textsc BLT} addresses this shortfall. Similarly, \cite{slagle2024spacebyte} corroborate the limitations of MegaByte and recommend enhancing the patching technique by dynamically segmenting on whitespaces and other whitespace-like bytes, as well as integrating a local encoder. Their findings reveal notable gains over tokenization-based transformer models under fixed computational budgets for specific domains, such as Github and arXiv, within the 1B parameter context. We also conduct experiments with this architecture, showing that although gains are possible, there remains a necessity for additional architectural innovation to ensure byte-level models can reliably scale and achieve parity with leading token-based models like Llama 3.

\section{Limitations and Future Work}

For this study, our architectural decisions are guided by the optimal training step counts established for Llama~3~\citep{dubey2024llama}. Nonetheless, it is important to note that these scaling laws were originally derived for BPE-level transformers and may not yield ideal ratios of data to parameter sizes when applied to {\textsc BLT}. Developing scaling laws specifically for {\textsc BLT} remains an open problem, which could, in turn, reveal even more favorable scaling behavior for our model. Furthermore, the bulk of our experiments have been limited to models with up to 1B parameters. Architectural optimizations might shift as we extend our investigations to 8B parameters and larger, potentially enabling additional performance gains at these larger scales.

Current transformer libraries and codebases are optimized primarily for architectures based on tokenizers. Although our work includes theoretically flop-equivalent experiments and leverages some optimized solutions—like FlexAttention—for layers that differ from standard transformer configurations, our implementation may still lag behind tokenizer-based models in wall-clock performance and could see improvements through additional optimization efforts.

Although {\textsc BLT} currently relies on an independently trained entropy model for patching, integrating the patching model into an end-to-end learning process presents a promising avenue for future research. In Section \ref{sec:distillation}, we introduce preliminary results suggesting that converting tokenizer-based models like Llama 3—those trained on upwards of 10T tokens—into byte-level models may be feasible. This is achieved by initializing the global transformer with pre-existing weights and then freezing those parameters. Continued exploration of this technique could reveal strategies that not only preserve the strengths associated with bytefying, but also potentially enable these models to outperform their tokenizer-based counterparts, all without necessitating complete retraining.

\section{Conclusion}
\label{section:conclusion}

In this work, we introduce the Byte Latent Transformer (\textbf{BLT}), an innovative architecture that departs from the traditional reliance on fixed-vocabulary tokenization in large language models. Through the implementation of a flexible, trainable technique for assembling bytes into patches, {\textsc BLT} dynamically distributes computational effort according to the underlying data’s intricacy. This approach yields notable advancements in both computational efficiency and resilience. Our comprehensive scaling experiments indicate that {\textsc BLT} is competitive with established tokenization-based architectures, such as Llama 3, reaching similar performance metrics at scales up to 8B parameters and 4T bytes. Notably, the architecture enables the reduction of inference flops by as much as 50\%, with only minimal compromise in downstream evaluation measures. Beyond these gains, {\textsc BLT} introduces the possibility of concurrent scaling of both the model and patch size, all while adhering to a fixed inference budget—an approach that proves advantageous for commonly encountered compute limitations in real-world scenarios. By directly processing raw byte sequences, {\textsc BLT} enhances the capacity to generalize across rare or noisy data, significantly improving robustness and providing richer representations of sub-word features. Taken together, these findings highlight {\textsc BLT} as a compelling and scalable replacement for conventional tokenization pipelines, equipping language models with greater adaptability and computational efficiency.

\end{document}